\documentclass[11pt,a4paper]{article}
\usepackage[margin=22mm]{geometry}
\usepackage{graphicx}
\usepackage[T1]{fontenc}
\usepackage{lmodern}
\usepackage{microtype}
\usepackage{booktabs}
\setlength{\parskip}{0.6em}
\setlength{\parindent}{0pt}
\renewcommand{\arraystretch}{1.15}

\title{\vspace{-14mm}\textbf{The other three agents}\\[2mm]\large Heuristics 1, Monte-Carlo, and Real-AI}
\author{}
\date{}

\begin{document}
\maketitle
\vspace{-12mm}

Companion to \texttt{h2-architecture.pdf}. Four agents can play this engine and
they disagree about what a decision \emph{is}: a weighted preference, an
estimate from simulated futures, a learned distribution, or a modelled change in
score. Each is drawn below as one decision, end to end.

\section*{Heuristics 1 --- \texttt{heuristic}}

\begin{center}\includegraphics[width=\textwidth]{h1.pdf}\end{center}

A phase table picks one evaluator per decision; the evaluator returns a numeric
weight per action; weights are normalised into a distribution and sampled.
Actions an evaluator declines fall to \texttt{defaultWeight}, scaled by the dice
odds when the action carries a target number.

Its weakness is the reason H2 exists: the weights are hand-tuned at the point of
use, so a number from the combat evaluator and a number from the site-phase
evaluator are not in the same unit and cannot be compared, checked, or
explained. It is nonetheless the strongest agent here --- H2 still loses to it.

\section*{Monte-Carlo --- \texttt{mc}}

\begin{center}\includegraphics[width=\textwidth]{mc.pdf}\end{center}

No tree, no network, no deck knowledge. It widens its own view into a playable
state, plays every candidate forward through the \emph{real engine} under a
random policy for a few turns, and keeps whichever candidate's playouts end with
the best average tournament-score differential.

Two departures from textbook rollouts, both about spending a small budget well:
\textbf{round-robin allocation} gives every candidate the same sample count, and
\textbf{common random numbers} compare all candidates against one determinized
world and one playout seed per round, so paired sampling removes between-candidate
noise that is far larger than the differences being measured.

The catch is stated in its own module: views the determinizer cannot handle ---
mid-chain, in combat, with pending effects --- are delegated to a fallback agent.
Those are a large share of all decisions, so \emph{the fallback materially shapes
how this agent plays}, which is the same trap H2 was in until the H1 fallback was
removed from it.

\section*{Real-AI --- \texttt{bc}}

\begin{center}\includegraphics[width=\textwidth]{bc.pdf}\end{center}

\textbf{This is a trained network, not a language model.} An action-conditioned
policy/value net is trained by \texttt{train/train\_bc.py} on recorded play; the
trainer exports weights as JSON, and the agent mirrors the forward pass in pure
TypeScript --- embeddings, a mean-pooled entity-set encoder, a state torso, a
per-candidate scorer with masked softmax, a tanh value head --- so no runtime or
native dependency is needed for a net this small.

Two details are worth knowing because they are where this design goes wrong
quietly:

\begin{itemize}\itemsep2pt
\item every weights file carries a \texttt{selfTest} block of real inputs and the
      trainer's outputs, replayed at load, and a card-vocabulary hash. The agent
      \emph{refuses to play} if either fails, so featurizer or weight drift
      cannot masquerade as a weak policy.
\item how the decision is read out of the distribution is a separate choice from
      the distribution. Argmax suits weights cloned from another agent's soft
      targets; weights cloned from human one-hots need class-mass or temperature
      sampling and \emph{score in single digits under argmax}.
\end{itemize}

\section*{What separates them}

\begin{center}\small
\begin{tabular}{lllll}
\toprule
 & \textbf{H1} & \textbf{MC} & \textbf{Real-AI} & \textbf{H2} \\
\midrule
a decision is & a weight & a simulated future & a learned distribution & a change in score \\
knowledge from & hand tuning & the engine & recorded play & rules + card data \\
unit & unitless & TSD & probability & TSD $\rightarrow$ win prob. \\
explains itself & no & by playout & no & rationale tree \\
needs a fallback & no & yes, often & no & no, since removed \\
cost per decision & negligible & rollouts & one forward pass & module evaluations \\
\bottomrule
\end{tabular}
\end{center}

Only two of the four can say why they did something. MC can show the playouts
that produced an average; H2 can print the rationale tree with every tunable
named. H1 and Real-AI both emit a number with no account of it --- which is
tolerable in a net, and was the specific defect H2 was built to remove from the
hand-written one.

\end{document}
